\section{The Spin Obstruction}

The committed substrate of this theory is the mesh $\{3,4,3,4\}$, whose docks carry $24$ sites built on the
$D_4$ root system. Those sites already hold a spinor. The honeycomb $\{5,3,4\}$ is a different object, and it is
worth a section of its own precisely because it is the nearest neighbor that fails. It is the order-$4$
dodecahedral honeycomb of three-dimensional hyperbolic space, and it sits one rung away from the committed mesh
along the family of regular hyperbolic honeycombs. Studying it is comparison and selection, not the thesis. It
shows why the committed mesh needs its spinor sites, and it shows what a substrate without one can still be
made to do.

The hardest fact about $\{5,3,4\}$ is short.

\begin{result}[The bare obstruction]
The bare knit on $\{5,3,4\}$ cannot make a fermion. The active sites of a dock on $\{5,3,4\}$ carry no spinor,
and no bound state, sum, or subrepresentation of the bare tone field will ever carry one.
\end{result}

This is not an engineering difficulty. It is a theorem. This section states the obstruction exactly, confirms
it by running the knit, and then shows that the \emph{geometry} of $\{5,3,4\}$ carries the spin structure even
though the \emph{knit} does not. That gap is closed in four ways, three that change the knit and one that
changes the substrate. The demonstrations of those routes, the experiments that show the fields actually move
and the defects actually carry spin, are the subject of the section on matter on $\{5,3,4\}$. The selection
argument this obstruction feeds, why spinor-carrying sites are rare across the whole family of honeycombs, is
the tessellation survey and the dimension ladder.

\begin{principle}
A spinor is projective. The bare tone field is linear. You cannot build the first from the second.
\end{principle}

\subsection{The substrate}

A short statement of the comparison substrate fixes the terms. Two words recur, so a plain definition of each
helps. A \emph{spinor} is the mathematical object that describes a particle of half-integer spin. Its signature
oddity is that it must be turned twice around, through $720$ degrees rather than $360$, before it returns to
itself. A \emph{fermion} is a particle of matter, an electron or a quark, described by a spinor. The question of
this section is whether the bare knit on $\{5,3,4\}$ can host such an object.

\begin{table}[ht]
\centering
\begin{tabular}{@{}ll@{}}
\toprule
property & value \\
\midrule
dimension & $3$ (hyperbolic, negatively curved) \\
dock shape & dodecahedron \\
rotational symmetry & icosahedral group $A_5$ \\
active sites & $12$ of the $24$ \\
the knit & reversible, charge-conserving lattice gas, collide then stream \\
\bottomrule
\end{tabular}
\caption{The comparison substrate $\{5,3,4\}$.}
\label{tab:534-substrate}
\end{table}

A dock holds $24$ sites, one tone each. The tone is ternary, a value in $\{-1, 0, +1\}$ read as pain, peace,
pleasure. The knit streams those tones along the sites of a dock, then collides them by a fixed reversible
charge-conserving table. On $\{5,3,4\}$ the knit moves the tone field along $12$ of the $24$ sites, the ones the
icosahedral group $A_5$ acts on. Those $12$ active sites are the whole stage the bare tone field lives on, and
the obstruction is a statement about what can transform on that stage.

\subsection{The obstruction, stated exactly}

The bare knit places the tone field on the $12$ active sites of each dock. The field is, at each dock, an
occupation of those $12$ directions.

Under the icosahedral group $A_5$ the $12$ sites \textbf{permute}. A permutation of basis vectors is a
\textbf{linear} representation. Its action is to shuffle the $12$ slots into one another, and shuffling slots is
the most basic linear thing a symmetry can do.

A spinor is a \textbf{projective} representation. It lives in the nontrivial class of the Schur multiplier
$\mathbb{Z}_2$. It is the representation that changes sign under a $2\pi$ rotation. The two cannot be bridged.

\begin{table}[ht]
\centering
\begin{tabular}{@{}lll@{}}
\toprule
object & representation class & sign under $2\pi$ \\
\midrule
the $12$-site occupation field & linear & $+1$ \\
a spinor & projective ($\mathbb{Z}_2$) & $-1$ \\
\bottomrule
\end{tabular}
\caption{The two representation classes that the obstruction separates.}
\label{tab:534-rep-class}
\end{table}

\begin{theorem}[The spin obstruction]
The $12$-site occupation field on $\{5,3,4\}$ is a linear representation of $A_5$. A spinor is a projective
representation in the nontrivial class of the Schur multiplier $\mathbb{Z}_2$. No tensor product, bound state,
sum, or subrepresentation of a linear representation is projective. Therefore the bare tone field carries no
spinor mode.
\end{theorem}

The proof is one line of cohomology. The cocycle of a tensor product \textbf{adds}. A linear representation has
cocycle \textbf{zero}. So every tensor product, every bound state, every sum, and every subrepresentation of the
bare field has cocycle zero \textbf{forever}. There is no hidden corner of the bare field where a spinor could
condense. The cocycle that a spinor needs is simply not in the algebra the bare field generates.

\begin{principle}
There is no emergent spinor mode hidden in the bare tone field. This is a theorem, not a difficulty.
\end{principle}

\subsection{The obstruction, confirmed dynamically}

A theorem about representation classes can be checked against the running knit, and it was. The lattice-gas
spinor experiment, at depth $L2$, rotates the $12$-site occupation field through a full $2\pi$ turn and reads
off the holonomy, the sign the field returns with.

\begin{table}[ht]
\centering
\begin{tabular}{@{}lll@{}}
\toprule
field & holonomy after $2\pi$ & verdict \\
\midrule
the bare $12$-site occupation & $+1$ & a \textbf{vector}, not a spinor \\
the spin bundle of the same geometry & $-1$ & a genuine spinor \\
\bottomrule
\end{tabular}
\caption{The dynamical confirmation. The occupation field returns to $+1$, a vector. The spin bundle of the same
geometry returns to $-1$, a spinor.}
\label{tab:534-holonomy}
\end{table}

The occupation field returns to $+1$. It is a vector. The spin bundle of the \textbf{same} $\{5,3,4\}$ geometry
returns to $-1$. So the spinor exists in the geometry but not in the field the bare knit lives on. That split is
the whole problem. The cocycle the knit needs is sitting in the geometry next door, untouched.

\subsection{The geometry carries the spin structure, at every scale}

Although the bare knit has no spinor, $\{5,3,4\}$ itself does, in the ultraviolet and in the infrared, at short
distance and at long. The frame bundle of the geometry realizes a genuine spinor, and four experiments, all at
depth $L2$ except where noted, establish it from independent directions.

\begin{table}[ht]
\centering
\begin{tabular}{@{}p{0.30\textwidth}p{0.10\textwidth}p{0.50\textwidth}@{}}
\toprule
experiment & depth & what it shows \\
\midrule
the cocycle & $L2$ & the icosahedral cover $2I$ is nonsplit, the cocycle is real, the split control $A_5 \times \mathbb{Z}_2$ fails \\
the spin connection & $L2$ & the Clifford lift of the $[5,3,4]$ Coxeter reflections gives every elementary edge loop $-1$ once around and $+1$ twice (short distance) \\
the continuum holonomy & $L2$ & by Gauss-Bonnet, a loop enclosing hyperbolic area $2\pi$ flips a spinor and leaves a vector blind (long distance) \\
spin and curvature & $L2$ & the genuine $2I$ spinor coexists with negative curvature (edge angle excess about $106$ degrees) \\
the icosian double cover & $L1$ & the icosahedral loop holonomy is the spinor double cover $2I$ \\
\bottomrule
\end{tabular}
\caption{The geometry of $\{5,3,4\}$ carries the spin structure at every scale.}
\label{tab:534-geometry-spin}
\end{table}

The cover is \textbf{genuine}, not a relabeled linear representation. The icosahedral cover $2I$ is perfect, so
its $-1$ is a commutator, and the split control $A_5 \times \mathbb{Z}_2$ cannot reproduce it. A relabeling could
have faked a sign, but a perfect group's commutator cannot be unwound that way, and the control confirms it.

The frame bundle realizes the spinor at \textbf{short} distance, where one elementary edge loop flips it, and at
\textbf{long} distance, where a Gauss-Bonnet loop enclosing hyperbolic area $2\pi$ flips it. A vector stays blind
in both cases. The spinor is not a short-distance artifact of the discretization, nor a long-distance artifact of
the continuum. It is present at both ends.

And the spinor sits on a \textbf{curved} bulk. This is the trade that flat $\{3,4,3,4\}$ cannot resolve from
curvature. The committed mesh has the spinor directions but is flat. The honeycomb $\{5,3,4\}$ has the curvature
but no field to carry the spinor. Each gift sits with the wrong partner.

\begin{result}[Spin present, field missing]
The spin structure of $\{5,3,4\}$ is present in the geometry, ultraviolet and infrared, on a curved bulk. The
only thing missing is a field that lives in it.
\end{result}

\subsection{The four routes that close it}

There are exactly four ways to put a field in that spin structure. Three change the knit. One changes the
substrate.

\begin{table}[ht]
\centering
\begin{tabular}{@{}p{0.22\textwidth}p{0.13\textwidth}p{0.32\textwidth}p{0.22\textwidth}@{}}
\toprule
route & changes & mechanism & cost \\
\midrule
1. defect & nothing & a disclination borrows the geometry's $-1$ holonomy & heavy soliton, rides on the unsolved persistence problem \\
2. spinor index and connection & the knit & stream a two-component spinor, transport it by an $\mathrm{SU}(2)$ connection with loop product $-1$ & adds a field \\
3. staggered Kahler-Dirac & the knit & put the fermion on the form complex, run $d + \delta$, then stagger to one Dirac field & adds the form structure \\
4. the $5$D $D_4$ pentacomb & the substrate & move to a honeycomb whose sites already carry the spinor & one extra dimension \\
\bottomrule
\end{tabular}
\caption{The four routes from the bare obstruction to a spinor field.}
\label{tab:534-four-routes}
\end{table}

\subsubsection{Route 1, the defect}

No knit change. The spinor is a topological disclination that carries the geometry's $-1$ holonomy directly. The
field does not carry the cocycle, the defect borrows it from the curved bulk. The cost is that the defect is a
heavy soliton, and it depends on solving the persistence problem first, since a defect that does not survive is
no use as matter.

\subsubsection{Route 2, add the field}

A knit change. Stream a two-component spinor and parallel-transport it by an $\mathrm{SU}(2)$ connection whose
loop product is the verified $-1$. The propagating-spinor experiment confirms it. The field returns minus itself
per $2\pi$ circuit and plus itself per two circuits, while the trivial-connection control returns plus. This is a
genuine propagating spinor, a field that moves around a loop and comes back with the spinor sign.

\subsubsection{Route 3, Kahler-Dirac}

A knit change. The form Dirac operator $d + \delta$ squares to the Hodge Laplacian on any oriented complex, which
the Kahler-Dirac experiment verifies. Stagger that operator to a single Dirac field. The forms are the spinor
bundle in disguise, so the fermion is recovered without ever leaving the differential structure the mesh already
carries.

\subsubsection{Route 4, the pentacomb}

No knit change. Move to the $5$D $D_4$ pentacomb $\{3,4,3,3,4\}$, whose $24$ sites are the $D_4$ roots. Those
roots split as $8v + 8s + 8c$, the vector and its two spinors, and the $8s$ and $8c$ carry the spinor $-1$
already. The bare directional knit, run on this substrate, has a spinor from the start.

\begin{principle}
Route $4$ is the only way to get a propagating spinor from the bare directional knit with no added field.
\end{principle}

The four routes are demonstrated in the section on matter on $\{5,3,4\}$, where the fields are shown to actually
move and the defects to actually carry spin. The pentacomb closure is the subject of the pentacomb resolution.

\subsection{Why the geometry is the right place to look}

The two gifts $\{5,3,4\}$ has that flat $\{3,4,3,4\}$ lacks are exactly the levers for its two hard problems, the
spin problem and the persistence problem.

\begin{table}[ht]
\centering
\begin{tabular}{@{}p{0.42\textwidth}p{0.50\textwidth}@{}}
\toprule
gift & consequence \\
\midrule
the spinor cover $2I$ is McKay-dual to $E_8$ & the natural home for rich matter and three generations \\
negative curvature gives a genuine bulk and boundary & hence holography, hence persistence by code \\
\bottomrule
\end{tabular}
\caption{The two structural gifts of $\{5,3,4\}$ and what each buys.}
\label{tab:534-gifts}
\end{table}

The spin problem is solved by moving to the spinor-carrying sites of the $D_4$ pentacomb, one extra dimension.
The persistence problem is solved by the holographic code on the boundary, treated in the section on holography.
The two assets and the two solutions line up. What $\{5,3,4\}$ cannot do from its bare sites, it does from $E_8$
and from the boundary.

\subsection{The one deep residual}

Most refinements are done. One question is genuinely open.

\begin{principle}
Can a single knit drawn from the eight atoms be \textbf{both} charge-conserving \textbf{and} a perfect tensor,
and likewise generate the spin connection dynamically?
\end{principle}

The causal structure is \textbf{free}. A bare reversible knit derives the causal wedge, as the holography section
shows. The perfect-tensor condition that the erasure code needs is the remaining constraint. That is the real
frontier of the comparison substrate.

\subsection{Index of verified results}

\begin{table}[ht]
\centering
\begin{tabular}{@{}p{0.74\textwidth}p{0.12\textwidth}@{}}
\toprule
claim & depth \\
\midrule
the bare knit gives no spinor mode (the obstruction, dynamical) & $L2$ \\
the spinor cover $2I$ is genuine (nonsplit) & $L2$ \\
the frame bundle carries spin, short distance & $L2$ \\
the frame bundle carries spin, long distance (Gauss-Bonnet) & $L2$ \\
spin and curvature coexist on $\{5,3,4\}$ & $L2$ \\
a propagating spinor from streaming plus the connection (route $2$) & $L2$ \\
the form Dirac operator squares to the Laplacian (route $3$) & $L2$ \\
the icosahedral loop holonomy is the double cover $2I$ & $L1$ \\
\bottomrule
\end{tabular}
\caption{The verified results behind the spin obstruction and its closures.}
\label{tab:534-verified}
\end{table}

The obstruction is exact, the geometry's spin structure is genuine and present at every scale, and the gap
between them is closed four ways. The comparison substrate fails to make a fermion from its bare sites, and the
shape of that failure is precisely the argument for the committed mesh, whose sites carry the spinor from the
start.
